% answers_short.tex —— 1986 年·广东卷·理工农医类 · 简短答案速查
% 由 tools/make_short_answers.py 自 answers.tex 的 \gksolhead/\gkanswer 提取，
% 并与 1986/1986_数学真题_广东理_答案要点.md 逐题比对后生成，勿手改（改源 MD 后重跑）。

\begin{gknotice}
\textbf{使用说明}：本节只给最终答案便于快速核对，完整推导见解析版。答案为编写者独立推导，\textbf{非官方标准答案}。
\end{gknotice}

\gkpart{大题}
\begin{itemize}[leftmargin=2.6em]
  \item[\textbf{1.}] (1) C；(2) B；(3) D；(4) D；(5) B；(6) E；(7) A；(8) $2156$；选项字母待核（若缺字方框为乘号，则为 E）；(9) $7200$；选项字母待核（若缺字方框为乘号，则为 C）；(10) B；(11) E；(12) A；(13) B；(14) C；(15) D（固定源末两项均印 D）；(16) B；(17) D；(18) A；(19) D；(20) 按“边”为线段理解时为 E；若指边所在直线，五项均不普遍成立，题意待核；(21) A；(22) C；(23) E；(24) C；(25) $f(x)=-\sqrt{2x-x^2}$，$0\le x\le1$；固定源答案表记录 E，但选项图缺失，字母无法核验
  \item[\textbf{2.}] (1) $y=(\sqrt3)^x$；(2) $x^2=y$（原题运算符位置呈方框字形；本答案按乘积解释，原字形待核）；(3) $\dfrac32$；(4) $4$；(5) $[-2,-1)$；(6) $180$；(7) $243$；(8) $\dfrac12$；(9) $\displaystyle x=k\pi+\frac{2\pi}{3}\text{ \text{或} }x=k\pi+\frac{5\pi}{6},\ k\in\mathbb Z$（按乘积解释；原题方框字形待核）；(10) $-2$
  \item[\textbf{3.}] (1) $x=2+\sqrt3$；(2) $a>0$，且唯一解为 $\displaystyle x=\frac{a+2+\sqrt{a^2+4a}}2$
  \item[\textbf{4.}] (1) 见解析版；(2) $\displaystyle \arcsin\frac{\sqrt{4b^2+\alpha^2}}{2\alpha}$
  \item[\textbf{5.}] (1) $\displaystyle a_n=2a_{n-1}\cos\alpha-a_{n-2}$；(2) 见解析正文；(3) $\displaystyle \frac{\cos((n+1)\alpha)\cos\frac{(2n+1)\alpha}{2}}{\cos\frac\alpha2}$
  \item[\textbf{6.}] $\displaystyle -\frac{2\sqrt{13}}{13}<m<\frac{2\sqrt{13}}{13}$
\end{itemize}
